% The Slide Definitions
%document
\documentclass[10pt]{beamer}
%theme
\usetheme{metropolis}
% packages
\usepackage{color}
\usepackage{listings}
\usepackage[ngerman]{babel}
\usepackage[utf8]{inputenc}
\usepackage{multicol}


% color definitions
\definecolor{mygreen}{rgb}{0,0.6,0}
\definecolor{mygray}{rgb}{0.5,0.5,0.5}
\definecolor{mymauve}{rgb}{0.58,0,0.82}

\lstset{
    backgroundcolor=\color{white},
    % choose the background color;
    % you must add \usepackage{color} or \usepackage{xcolor}
    basicstyle=\footnotesize\ttfamily,
    % the size of the fonts that are used for the code
    breakatwhitespace=false,
    % sets if automatic breaks should only happen at whitespace
    breaklines=true,                 % sets automatic line breaking
    captionpos=b,                    % sets the caption-position to bottom
    commentstyle=\color{mygreen},    % comment style
    % deletekeywords={...},
    % if you want to delete keywords from the given language
    extendedchars=true,
    % lets you use non-ASCII characters;
    % for 8-bits encodings only, does not work with UTF-8
    frame=single,                    % adds a frame around the code
    keepspaces=true,
    % keeps spaces in text,
    % useful for keeping indentation of code
    % (possibly needs columns=flexible)
    keywordstyle=\color{blue},       % keyword style
    % morekeywords={*,...},
    % if you want to add more keywords to the set
    numbers=left,
    % where to put the line-numbers; possible values are (none, left, right)
    numbersep=5pt,
    % how far the line-numbers are from the code
    numberstyle=\tiny\color{mygray},
    % the style that is used for the line-numbers
    rulecolor=\color{black},
    % if not set, the frame-color may be changed on line-breaks
    % within not-black text (e.g. comments (green here))
    stepnumber=1,
    % the step between two line-numbers.
    % If it's 1, each line will be numbered
    stringstyle=\color{mymauve},     % string literal style
    tabsize=4,                       % sets default tabsize to 4 spaces
    % show the filename of files included with \lstinputlisting;
    % also try caption instead of title
    language = Python,
	showspaces = false,
	showtabs = false,
	showstringspaces = false,
	escapechar = ,
}

\def\ContinueLineNumber{\lstset{firstnumber=last}}
\def\StartLineAt#1{\lstset{firstnumber=#1}}
\let\numberLineAt\StartLineAt



\newcommand{\codeline}[1]{
	\alert{\texttt{#1}}
}


% Author and Course information
% This Document contains the information about this course.

% Authors of the slides
\author{Felix Döring, Felix Wittwer, Anton Obersteiner}

% Name of the Course
\institute{Python-Kurs}

% Fancy Logo 
\titlegraphic{\hfill\includegraphics[height=1.25cm]{../templates/fsr_logo_cropped}}



% Custom Bindings
% \newcommand{\codeline}[1]{
%	\alert{\texttt{#1}}
%}

\usepackage{tikz}

% Presentation title
\title{Iterierbare Datenstrukturen}
\date{\today}


\begin{document}

\maketitle

\begin{frame}{Gliederung}
	\setbeamertemplate{section in toc}[sections numbered]
	\tableofcontents
\end{frame}


% #############################################################################
% -------------------- Wiederholung/Abschluss Objekte -------------------------
% #############################################################################
\section{Wiederholung/Abschluss Objekte}

\begin{frame}{Objekte und Vererbung}
	\begin{itemize}
		\item Was sind Klassen, was Objekte, was ist Vererbung?
		\item<2-> Was braucht unser Planet noch? {\tt class Planet(Vector):}
	\end{itemize}
	\begin{tikzpicture}
		\onslide<2->{
			\node         (origin)  at (0,   0) {$O$};
			\node         (planet)  at (5,   2) {};
			\node[above]  (Planet)  at (5,   2) {\tt planet};
			\draw[-latex] (origin) -> (planet);
		}
		\onslide<3->{
			\node         (heading) at (6.5, 1) {};
			\node[right]  (Heading) at (6.5, 1) {\tt vel};
			\node[above]  (mass)    at (5,   3) {\tt mass};

			\draw[-latex] (planet) -> (heading);
			\draw (planet) circle[radius=1];
		}
	\end{tikzpicture}
\end{frame}

\begin{frame}{Planeten-Eigenschaften}
	\lstinputlisting{resources/02_grundlagen/planet_question.py}
	\url{github.com/AntonObersteiner/python-lessons/tree/master/latex/slides/resources/02_grundlagen/planet_question.py}
\end{frame}
\begin{frame}{Planeten-Eigenschaften}
	\begin{tikzpicture}
		\node         (origin)   at (0,   0) {$O$};
		\node         (planet1)  at (5,   2) {};
		\node[right]  (Planet1)  at (5,   2) {\tt self};
		\node         (heading1) at (3.5, 3) {};
		\node[above]  (Heading1) at (3.5, 3) {\tt self.vel};
		\node[below]  (mass1)    at (5,   1) {\tt self.mass};

		\node         (planet2)  at (-2,  1) {};
		\node[left]   (Planet2)  at (-2.7,1) {\tt other};
		\node         (heading2) at (-2,2.5) {};
		\node[left]   (Heading2) at (-2,2.5) {\tt other.vel};
		\node[below]  (mass2)    at (-2, .3) {\tt other.mass};

		\node         (force1)   at ({5 - .35 * 7},   {2 - .35 * 1}) {};
		\node  (headingforce1)   at ({3.5-.35 * 3.5}, {3 - .35 *.5}) {};

		\draw[-latex] (origin) -> (planet1);
		\draw[-latex] (origin) -> (planet2);
		\draw[-latex] (planet1) -> (heading1);
		\draw[-latex] (planet2) -> (heading2);
		\draw (planet1) circle[radius=1];
		\draw (planet2) circle[radius=.7];
		\onslide<5-> \draw[-latex, red] (planet1) -> node[sloped, above] {\tt diff} (planet2);
		\onslide<7-8> \draw[blue] (0, 2.5) -- node[sloped, above] {\tt F\_abs} (2.5, 2.5);
		\onslide<8-> \draw[-latex,blue] (planet1) -> node[sloped, above] {\tt force} (force1);
		\onslide<10-> \draw[-latex,blue] (heading1) -> node[sloped, below] {$\mathrm{d}v$} (headingforce1);
	\end{tikzpicture}
	$$\begin{aligned}
		\dot s &= v && \onslide<2->{\text{{\tt self += vel * dt}\phantom M \alert{Planet(Vector)}}} \\
		\dot v &= a && \onslide<3->{\text{\tt vel  += acc * dt}} \\
		F &= m\cdot a && \onslide<4->{\text{\tt acc = force / mass}} \\
		d &= s_2 - s_1 && \onslide<5->{\text{\tt diff = other - self}} \\ % \phantom M \color{blue} (Vector.\_\_sub\_\_(...))} \\
		|F_G| &= \onslide<6->{
			G\cdot\frac{m_1 \cdot m_2}{|d|^2}
		} && \onslide<7->{
			\text{\tt F\_abs = G * \textcolor{blue}{s}.mass * \textcolor{blue}{o}.mass / dist ** 2}
		} \\
		F_G &= \onslide<8->{
			\frac{d}{|d|} \cdot |F_G|
		} && \onslide<9->{
			\text{\tt force = diff * (F\_abs / diff.abs())}
		} \\
	\end{aligned}$$
\end{frame}

% #############################################################################
% ----------------------------------- Listen ----------------------------------
% #############################################################################
\section{Listen}

\begin{frame}{list - und Schleifen}
	\lstinputlisting{resources/03_iterables/list_short.py}
	\begin{itemize}
		\item enthält variable Anzahl von Objekten mit fester Reihenfolge
		\item kann verschiedene Datentypen enthalten
		\item Auch Listen können in Listen gespeichert werden!
		\item Speicherverständnis!
	\end{itemize}
	\onslide<2->\lstinputlisting{resources/03_iterables/list_copy.py}
\end{frame}
\begin{frame}{list - Unpacking-Beispiel}
	\lstinputlisting{resources/03_iterables/list_unpacking.py}
	\begin{itemize}
		\item {\tt tuple} enthalten mehrere Elemente,
		\item sind aber unveränderbar $\to$ Vor- und Nachteile
	\end{itemize}
	\onslide<2->\lstinputlisting{resources/03_iterables/tuple_return.py}
\end{frame}
\begin{frame}{list - Anhängen und map}
	\onslide<1->\lstinputlisting{resources/03_iterables/list.py}
	\onslide<2->\lstinputlisting{resources/03_iterables/list_roots.py}
\end{frame}
\begin{frame}{list - filter}
	Primzahlen mit {\tt filter(func: data -> bool, data)} finden
	\onslide<2->\lstinputlisting{resources/03_iterables/list_filter.py}
\end{frame}


% #############################################################################
% --------------------------------- Dict, Set ---------------------------------
% #############################################################################
\section{Dict, Set}

\begin{frame}{dict}
	Zuordnung: $key \to value$
	\lstinputlisting[lastline=11]{resources/03_iterables/dict.py}
\end{frame}


\begin{frame}{Es gibt Mengen in Python}
	\begin{itemize}
		\item kann nur hashbare Einträge enthalten
		\item {\tt students = \{ "{}Elia", "Muzaina", ... \}}
		\item Fun fact: Python kann Unicode, aber mein \LaTeX{} nicht
	\end{itemize}
\end{frame}

% #############################################################################
% --------------------------------- Iterations --------------------------------
% #############################################################################
\section{Iterations}

\begin{frame}{Iteration}
	\begin{itemize}
		\item {\tt for $a$ in $A$: $do$}
		\item $a$ wird hier $A$ einmal durchlaufen,
		\item für jedes $a$ wird der block $do$ ausgeführt
		\item für Iterationen über Integer: \texttt{range([start], stop, step=1)}
		\item \alert{\texttt{\_\_iter\_\_}} Methode wird intern aufgerufen
		\item<2-> Gegensatz: {\tt while} mit Bedingung
	\end{itemize}
\end{frame}


% #############################################################################
% ------------------------------- File Handling -------------------------------
% #############################################################################
\section{File Handling}

\begin{frame}{File Handling}
	\begin{itemize}
		\item Dateien können mit \alert{\texttt{open(filename, mode="r")}} geöffnet werden
		\item \textit{File Handler} sind Iteratoren über die Zeilen einer Datei
		\item \textbf{Wichtig:} File Handler müssen auch wieder geschlossen werden
		\item \texttt{r} steht für Lesezugriff,  \texttt{w} für Schreibzugriff
	\end{itemize}
	\textbf{Beachte:} Wird eine Datei mit Schreibzugriff geöffnet, wird sie geleert! Also wichtige Inhalte vorher auslesen.
\end{frame}

\begin{frame}{File Handling - Beispiel}
	\lstinputlisting{resources/03_iterables/file.py}
\end{frame}


% #############################################################################
% ------------------------------- Import --------------------------------------
% #############################################################################
\section{Import}

\begin{frame}{Import}
	\begin{itemize}
		\item<1-> Funktionen, Klassen, ... aus Bibliotheken importieren
		\item<1-> {\tt import turtle} $\to$ {\tt turtle.goto(100, 100)}
		\item<2-> {\tt from math import sin} $\to$ {\tt sin(angle)}
		\item<3-> {\tt from math import *} $\to$ {\tt sin, cos, pi, exp, ...}
		\item<4-> {\tt from matplotlib import pyplot as plt}
		\item[]<5-> $\to$ {\tt plt.plot(X, Y); plt.show()}
		\item<6-> {\tt from Planet import Planet}
		\item[]<7-> $\to$ import {\tt Planet} from file {\tt Planet.py}
	\end{itemize}
\end{frame}


\end{document}
